\section{Results}
\label{sec:exp}

We first establish what compression does to the capability vector on
the heterogeneous panel (\S\ref{sec:phenomena}), then test prediction
arm by arm with one template (inputs and baselines; source axis;
configuration axis; independent tests and boundaries) for pruning
(\S\ref{sec:prune}), quantization (\S\ref{sec:quant}), and distillation
(\S\ref{sec:distill}), and close with a cross-arm synthesis and the
limitations (\S\ref{sec:synthesis}). Mechanism evidence, recovery
training, and the re-analysis of an earlier ability-space grid are
supporting material in Appendix~\ref{sec:app_support}.

\subsection{Capability responses on the heterogeneous panel}
\label{sec:phenomena}

\input{tables/panel_prune}

\paragraph{Pruning.} Table~\ref{tab:panel_prune} gives the
per-capability damage at $d=0.7$ and $d=0.5$, the density at which
math damage first exceeds one nat, and how QA ranks among the three
capabilities before that cliff. Three things are visible. Magnitudes
differ by two orders across series at the same density (OLMo-3-32B
$+0.01$ against Gemma-3-27B $+4.0$ on math at $d=0.7$). The ordering of
capabilities is series-specific: in the three Qwen3 sizes and in
OLMo-3-32B, QA is the least-damaged capability at every pre-cliff
density and improves, by up to $-0.79$ nats at Qwen3-4B; in every
Gemma-3 size, in Gemma-4-31B, and in Muse-30B, QA is the most-damaged
capability at the mildest densities (Gemma-4-31B: $+2.89$ QA against
$+0.68$ math at $d=0.7$). OLMo-3-7B falls in between with differences
of at most $0.15$ nats. A noise screen sets the scale of what counts:
the median $|\Delta L|$ at $d=0.9$ is $0.013$ nats, and of the 62
cells with a nominal improvement across all arms, 29 exceed twice that
median; the rest are treated as zero. Cliffs occur between $d=0.5$ and
$0.7$ for Qwen3 and Gemma, at $0.55$ for Muse-30B, at $0.35$ for
OLMo-3-7B, and had not occurred by $d=0.3$ for OLMo-3-32B, whose QA
loss is $-0.24$ nats below dense at that density. Robustness is not
monotone in size within a series: Gemma-3-12B is the most robust
Gemma-3 size and 27B collapses earlier. We read all of this as
heterogeneity that a predictor has to reproduce, not as an invariant.

\paragraph{Quantization.} Appendix Table~\ref{tab:panel_quant} covers 14
models. Int8 and int6 stay within $0.01$ nats of dense throughout.
Int4 damages small models (Gemma-3-270M $+2.0$, $+2.4$, $+0.9$) and
barely touches large ones (Gemma-3-27B $+0.33$, $+0.23$, $-0.67$), and
the QA improvement under int4 recurs in five models (Qwen3-1.7B, 8B,
14B, Gemma-3-12B, Gemma-3-27B), all clearing the noise screen. Int3
collapses 12 of the 14 models by $4.6$ to $28$ nats; OLMo-3-7B loses
$2.8$ to $4.7$ nats and OLMo-3-32B only $0.4$ to $1.2$, so the
collapse is common but not universal, and OLMo-3 is the robust series
under both interventions. Within the tested bit-widths, a learned
exponential $q_c e^{-k b}$ fits held-out bits better than the fixed
$q_c 4^{-b}$ shape (pooled MAE $6.7$ against $7.6$ nats), which shows
that the fixed base is wrong for this quantizer and nothing more; the
squared step-size ratio of the quantizer between 4 and 5 bits is
$(7/15)^2\approx0.22$ (Appendix~\ref{sec:app_compress}), and a loss
ratio would follow from it only under additional error-distribution
assumptions that we do not test.

\paragraph{Distillation.} Black-box traces from two API teachers (600
per domain each) train Gemma-3 students under a full-trace and an
answer-only recipe in a paired design. The answer-only recipe hurts
math more than full-trace training ($+0.19$ to $+0.22$ against $+0.09$
to $+0.12$ nats relative to the initial student), while QA improves
under every recipe ($-0.58$ to $-1.44$). Because the recipes differ in
supervised tokens and update counts, this pairing does not isolate
reasoning coverage; the controlled comparisons that hold tokens fixed
are those of \S\ref{sec:distill}.

\paragraph{What the endpoint does and does not say.} The same
capability heterogeneity appears on independent secondary benchmarks
for math and QA, where the primary-benchmark response predicts the
secondary one better than any other capability's response does, with
bootstrap intervals excluding zero in most hold-out protocols; for code
the cross-capability predictor is as good or better, so code
specificity is not established at this resolution
(Appendix~\ref{sec:app_probes}). A lower QA loss does not by itself
mean better QA behavior: in the four pruned OLMo-3-32B cells whose QA
loss lies below dense, exact-match accuracy does not rise
(Appendix~\ref{sec:app_probes}). The phrase ``QA improves'' in this
paper refers to the loss.

\subsection{Pruning arm}
\label{sec:prune}

\input{tables/pred_main}

\paragraph{Inputs, candidates, baselines.} Inputs are $N_0$, $D_0$,
and the dense capability loss $L_{0,c}$ of the source; the candidates
and baselines are those of \S\ref{sec:predictors}. All fits use the
nine Pythia development states.

\paragraph{Source axis (Table~\ref{tab:pred_main}; Appendix Table~\ref{tab:pred_source}).} With the
density held fixed and one pretraining stage held out, the
configuration-indicator regression reaches MAE $0.253$ (math), $0.292$
(code), and $0.806$ (QA) against $0.413$, $0.516$, and $0.500$ for the
strongest simple baseline; the improvement interval excludes zero for
code only, and for QA the regression is worse than predicting zero
change. Same-input comparisons say which input carries the gain: adding
the dense anchor $L_{0,c}$ to $\{N_0,D_0\}$ has intervals above zero
for math ($+0.55$, $[0.29,1.00]$) and code ($+0.66$, $[0.38,1.13]$),
whereas adding $D_0$ to $\{N_0,L_{0,c}\}$ does not ($+0.36$,
$[-0.04,0.60]$; $+0.40$, $[-0.01,0.63]$); on an earlier split that also
held out sizes, the $D_0$ interval cleared zero in one pruning cell
(math) out of six. On a frozen prospective at a new stage (step 96k on
the 160M and 1.4B sources, between 64k and 143k) the frozen regression
beat both its no-$D_0$ variant and the per-configuration median on all
three capabilities, with the largest margin on math.

\paragraph{Configuration axis (Table~\ref{tab:pred_main}; Appendix Table~\ref{tab:pred_config_prune}).}
On the development sources plus a new source (410M at 96k), pooled over
the unseen densities $0.65$ and $0.55$, the frozen shared power form
has MAE $0.371$, $0.524$, $0.759$ against $1.37$, $1.64$, $1.61$ for
the strength-only curve: a source-blind curve cannot see that one state
(160M at 143k) is several times more fragile than another. The
same-input regression A2, added after the test, is better on math and
code and within noise on QA, so the power structure adds nothing that
the per-density regression lacks. Two frozen rounds on new sizes then
set the range. For the 1B source, inside the development size range, at
the interpolated densities over two stages the power form, A2, and the
source-free median curve are within $0.05$ nats of one another on every
capability; per stage (Appendix~\ref{sec:app_p1v2}) A2 is best at 112k
on each capability and the median curve is best at 32k, driven by QA.
For the 6.9B source, five times the largest development size, the
source-conditioned forms lose to the median curve on math and QA and
tie on code: the regression extends the panel's size trend into QA
improvements of one to three nats that do not occur. At density $0.55$
every source-conditioned form fails on every source (errors $0.6$ to
$4.5$ nats) and a source-free curve is closest; the single-stage 1B@96k
prospective (Appendix Table~\ref{tab:pred_config_full}) shows the same
two facts. Under protocol B (early stages only; Appendix
Table~\ref{tab:p1v2}) the 6.9B interpolation errors shrink to the level
of the source-free curves without beating them.

\paragraph{Boundaries.} Source-state information improves pruning
prediction at seen densities and at interpolated densities, for
sources inside the development size and stage range, and mostly for
math and code. The gain does not require the shared power structure.
It does not extend to a source five times larger, to density $0.55$, or
reliably to QA, whose response the current inputs do not carry.

\subsection{Quantization arm}
\label{sec:quant}

\paragraph{Inputs, candidates, baselines.} The same inputs; the
configuration-indicator regression per bit-width against its no-$D_0$
variant, the per-bit mean and median, and zero change.

\paragraph{Source axis.} Pooled over the four bit-widths with a stage
held out (Table~\ref{tab:pred_main}), the regression improves on the
per-bit median by $+0.27$ (math), $+0.93$ (code, interval excluding
zero), and $-0.52$ (QA); adding $D_0$ has an interval above zero for
code ($+1.45$, $[0.15,2.10]$), a positive point estimate for math, and
none for QA. The pooled number is dominated by int3, so the frozen
step-96k test on three new sources is reported by regime: at 4 bits and
above the responses are near zero and candidate and median errors are
both small ($0.03$ against $0.08$ for math, $0.10$ against $0.09$ for
code, $0.17$ against $0.13$ for QA), which demonstrates no gain at that
error scale and does not show that none exists; at int3 the candidate
is better for math ($+0.45$) and worse for code ($-1.12$) and QA
($-0.32$), so the magnitude of the collapse does not transfer
consistently across capabilities.

\paragraph{Configuration axis (Table~\ref{tab:pred_main}; Appendix Table~\ref{tab:pred_config_qd}).} The
tested settings are the integer bit-widths; there is no integer between
3 and 4, so the unseen-configuration tests are the 5-bit rule and the
unseen sizes. For the in-range 1B source the regression is worse than
the per-bit median at 4 bits and above, and at int3 it is better for
math and code ($+1.2$ nats each) and worse for QA ($-1.0$). For the
6.9B source it is worse everywhere, by $0.1$ to $0.4$ nats above 3 bits
and by $1.2$ to $3.6$ nats at int3, where it overshoots the collapse.
Under protocol B (Appendix Table~\ref{tab:p1v2}) the 6.9B differences
above 3 bits shrink to at most $0.05$ nats and at int3 the regression
is ahead at 112k ($1.29$ against $3.8$ to $4.1$) and behind at 32k. The
single-stage 1B@96k test had the regression ahead on the int3 aggregate
($0.40$ against $0.62$ for the median and $2.46$ without $D_0$) with
the same math-positive, code-negative pattern per capability.

\paragraph{Boundaries.} In the near-zero regime (4 bits and above) no
predictor is distinguishable from a per-bit constant on these panels.
In the collapse regime (int3) the source-conditioned regression carries
information for in-range sizes, with a sign that differs by capability,
and it does not transfer to a five-times larger source. Which of these
statements holds depends on the fitting protocol, as the protocol B
exceptions show, so a quantization result in this paper is always
stated with its precision regime and protocol.

\subsection{Distillation arm}
\label{sec:distill}

\paragraph{Inputs, candidates, baselines.} The response is the signed
transfer $\delta_c$ relative to the initial student. On the source
axis the inputs are the student's $N_0$, $D_0$, and $L_{0,c}$ with a
fixed teacher and pool; on the pool axis the inputs are the processed
tokens $T$ and the reuse count $E=T/D_U$ of a fixed student.

\paragraph{Source axis (Table~\ref{tab:pred_main}).} With one stage
held out on the nine-state LoRA panel, the linear model in
$\{N_0,L_{0,c},D_0\}$ halves the error of a constant for math ($0.014$
against $0.031$), ties for code ($0.057$ against $0.061$, with the
no-$D_0$ variant at $0.044$), and is worse than a constant for QA
($0.263$ against $0.159$). Nine cells give no usable interval; these
are point estimates.

\paragraph{Configuration axis (Table~\ref{tab:pred_main}; Appendix Table~\ref{tab:pred_config_qd}).} Four
forms frozen on the endpoint pools U75 and U600 were applied to an
unseen middle pool U225, whose reuse count ($2.3$ to $4.7$) lies inside
the endpoint range; because the endpoints are pool prefixes and U225 is
randomly sampled, the sampling protocol changes as well. The reuse-count
form beats the constant only for QA ($0.706$ against $1.449$) and ties
the joint form; for math and code the constant is best ($0.122$,
$0.069$) and the fitted forms over-extrapolate. All four forms were
frozen and are reported; the per-capability recommendation (reuse
count for QA, constant for math and code) is a reading made after the
result. The residual is dominated by systematic bias: the QA reuse-count
prediction is $-0.87$ nats off at one budget and $+0.48$ at the other,
a form misfit, and around that bias the variation across pool samples
exceeds the variation across training seeds (QA $0.43$ against $0.04$
nats), which describes this design of three pools and two seeds and
decomposes no population variance.

\paragraph{Boundaries.} A transfer predictor from the student state
exists for math on the controlled panel; for QA a constant is better on
both axes tested; the reuse-count coordinate carries information for QA
on the unseen pool with a budget-dependent bias. Both panels are small,
the pool test changes two things at once, and the absolute-exposure
matrix that separates them is outside the data cut-off of this paper.

\FloatBarrier
\subsection{Cross-arm synthesis and limitations}
\label{sec:synthesis}
\label{sec:planned}

\paragraph{Which inputs help, and when.} The dense capability vector
of the source is the input with the clearest support: on the
controlled pruning panel adding it has intervals above zero for math
and code, and every successful frozen test uses it. Training tokens
$D_0$ add information beyond $\{N_0,L_{0,c}\}$ with interval support in
one cell (quantization, code) and positive point estimates in most
others; they never help QA. Source information beats source-free
curves for pruning inside the development range and at seen or
interpolated densities, and for quantization only in the collapse
regime and for in-range sizes.

\paragraph{Which forms add nothing.} On the same inputs, the shared
power form never beats the per-density regression with a fixed
interpolation rule, so no exponent is established by these data. At a
new in-range size the per-density regression matches the source-free
median curve on the pooled stages and beats it at the late stage only;
at a five-times larger size it loses.

\paragraph{Which transfers were tested.} Unseen stage inside the
range: supported for pruning (leave-one-out and frozen). Unseen density
inside the fitted range: supported for in-range sizes. Density $0.55$:
fails for every source-conditioned form. Unseen size inside the range:
on par with source-free curves. Size five times beyond the range: fails
in both arms. Unseen bit-width: only the 5-bit rule, which is worse
than a per-bit constant. Unseen pool: reuse count helps QA only. Each
of these is one or two source states or one pool; none is a population
statement.

\paragraph{Resource-ratio models.} A method-agnostic response in the
storage or active-parameter ratio was tested on the heterogeneous
panel: at the common storage ratio $0.5$, pruning to $d=0.5$ costs
median $+5.9$, $+8.2$, and $+2.9$ nats while int8 quantization costs
none; over 75 storage-matched pairs (14 with both sides pre-cliff) the
mean absolute method gap is $9.3$ nats with sign agreement $0.61$, and
the two shared-law hypotheses fit on pre-cliff cells reach held-out
sign accuracies of $0.51$ and $0.37$. These particular models are
rejected; whether another single form could unify the arms is untested.

\paragraph{From predictors to selection.} A capability- and cost-aware
method selection would use predictors like these as loss constraints.
An offline replay with frozen predictions exists and quantifies the
decision cost of prediction error, but no held-out agent or selection
validation was run, so nothing here is a validated selector.

\paragraph{Limitations.} The controlled panels have three sizes,
three stages, and three pools, so intervals are conditional on the
panel. QA is the capability the current inputs carry worst on every
axis. The quantization conclusions depend on the precision regime and
on the fitting protocol. The distillation pool test changes sampling
protocol and pool size together. Code specificity of the endpoint was
not detected at the available resolution. Recovery training, geometry,
and the M1/M2 shared-law tests rest on single models or pre-cliff
subsets and are reported as supporting evidence only.
